% !TEX root = ../thesis.tex

%=========================================================
% \setcounter{section}{-1}
\section{Quasiparticles and Transport Regimes}
\label{sec:quasi}
%=========================================================
    
    This section establishes the transport framework used throughout the chapter. We begin with a many-electron Hamiltonian and reduce it to the quasiparticle parameters that govern low-energy transport. In a large diffusive conductor these parameters lead to the Drude description and the measured normal-state resistance. When phase coherence extends across a finite conductor, scattering theory resolves transport into transmission channels. A coherent diffusive conductor contains many such channels, whereas an atomic contact is described by a small channel set.

    The aim is not to reproduce each theory in full, but to establish the normal-state parameters and length-scale hierarchy that select the applicable transport description. Bulk diffusive transport is described by Drude theory, while coherent finite conductors require a channel-based scattering description. \cite{ashcroft_solid_2012, gross_festkorperphysik_2014}

    %=========================================================
    \subsection{Microscopic Quasiparticles}
    \label{subsec:quasi:micro}
    %=========================================================

        We consider $N$ electrons at coordinates $\{\vec{r}_i\}$ in a crystalline solid whose ionic cores are fixed at positions $\{\vec{R}_a\}$. With the ionic coordinates absorbed into $V_\text{ion}$, the Hamiltonian in first quantization is
        \begin{equation}
            \hat{H}
            = \sum_{i=1}^N \left(
                \frac{\hat{\vec{p}}^{\,2}_i}{2m}
                + V_\text{ion}(\vec{r}_i)
            \right)
            + \frac{1}{2} \sum_{i \neq j}
            \frac{e^2}{4\pi\varepsilon_0 \lvert \vec{r}_i - \vec{r}_j \rvert}\,,
            \label{eq:quasi:full-hamiltonian}
        \end{equation}
        with canonical momenta $\hat{\vec{p}}_i = -\ima\hbar\nabla_i$ and electron mass $m$. The ionic potential $V_\text{ion}(\vec{r})$ represents the Coulomb attraction of the electrons to the positively charged lattice. To a good approximation it is periodic,
        \begin{equation}
            V_\text{ion}(\vec{r} + \vec{R})
            = V_\text{ion}(\vec{r})
            \label{eq:quasi:ionic-periodicity}
        \end{equation}
        for all Bravais lattice vectors $\vec{R}$.

        The many-electron wave function $\Psi(\vec{r}_1,\dots,\vec{r}_N,t)$ obeys the time-dependent Schrödinger equation
        \begin{equation}
            \ima\hbar \frac{\partial}{\partial t} \Psi
            = \hat{H}\,\Psi\,,
            \label{eq:quasi:schroedinger}
        \end{equation}
        supplemented by antisymmetry under particle exchange due to Fermi statistics. Solving this problem exactly is not feasible. To arrive at a tractable description of electronic transport we now introduce a sequence of simplifications. \cite{ashcroft_solid_2012, gross_festkorperphysik_2014}

        %=========================================================
        \subsubsection*{Effective One-Electron Description}
        %=========================================================

            First, according to the Born--Oppenheimer approximation, the fast electronic motion is separated from the slow ionic motion, and the ions are treated as fixed sources of the ionic potential. For an ideal lattice this potential is periodic, while impurities and structural defects are represented separately as disorder. Next, electron--electron interactions are taken into account in an average way. Screening in a metal strongly reduces the long-range Coulomb repulsion, and the dominant effect of interactions can be encoded in an effective potential and a renormalized band structure.
            On this level, the complicated many-electron problem is reduced to an effective one-electron eigenvalue problem
            \begin{equation}
                \hat{H}_\text{eff}\,\psi(\vec{r})
                = E\,\psi(\vec{r})\,,
                \qquad
                \hat{H}_\text{eff}
                = \frac{\hat{\vec{p}}^{\,2}}{2m} + V_\text{per}(\vec{r})\,,
                \label{eq:quasi:heff}
            \end{equation}
            where $V_\text{per}(\vec{r})$ is a periodic effective potential that already includes the mean-field action of the other electrons. 
            Residual interactions and disorder manifest at low energies through renormalized band parameters and relaxation processes that we summarize below by a small set of effective quasiparticle parameters.

            The stationary eigenvalue problem with a periodic potential,
            \begin{equation}
                \left[
                    -\frac{\hbar^2}{2m}\nabla^2 + V_\text{per}(\vec{r})
                \right] \psi(\vec{r})
                = E\,\psi(\vec{r})\,,
                \label{eq:quasi:schroedinger-periodic}
            \end{equation}
            is the starting point for band theory. Its solutions provide the electronic states from which the low-energy transport parameters are constructed. \cite{born_zur_1927, ashcroft_solid_2012,bloch_uber_1929}

        %=========================================================
        \subsubsection*{Bloch States and Band Structure}
        %=========================================================

            Because $V_\text{per}(\vec{r})$ has the periodicity of the lattice, Bloch's theorem applies. The eigenstates of $\hat{H}_\text{eff}$ can be chosen as Bloch functions
            \begin{equation}
                \psi_{n\vec{k}}(\vec{r})
                = \mathrm{e}^{\ima\vec{k}\cdot\vec{r}}\,
                u_{n\vec{k}}(\vec{r})\,,
                \label{eq:quasi:bloch}
            \end{equation}
            where $n$ is a band index, $\vec{k}$ lies in the first Brillouin zone, and $u_{n\vec{k}}(\vec{r}+\vec{R}) = u_{n\vec{k}}(\vec{r})$ is lattice-periodic. Inserting Eq.~\eqref{eq:quasi:bloch} into Eq.~\eqref{eq:quasi:schroedinger-periodic} yields for each $\vec{k}$ a discrete set of eigenvalues $E_n(\vec{k})$ that define the band structure.

            The Bloch functions form an orthonormal basis. In the independent-electron approximation, the many-electron ground state is obtained by filling the one-particle eigenstates up to the Fermi energy $E_\mathrm{F}$, with at most two electrons, spin up and spin down, per state. For a given electron density, the filling of the bands fixes $E_\mathrm{F}$ and the corresponding Fermi surface defined by
            \begin{equation}
                E_n(\vec{k}) = E_\mathrm{F}\,.
                \label{eq:quasi:fermi-surface}
            \end{equation}

            In simple metals one or a few bands cross the Fermi level. For transport at low temperature and low bias, only states in a narrow energy window $\lvert E - E_\mathrm{F} \rvert \sim \max(k_\mathrm{B}T, eV)$ are relevant. The detailed band structure away from the Fermi surface can therefore be compressed into effective parameters such as the Fermi velocity and the density of states at the Fermi energy. These are the quantities that enter the transport models below. \cite{ashcroft_solid_2012,bloch_uber_1929}

        %=========================================================
        \subsubsection*{Semiclassical Dynamics and Scattering}
        %=========================================================

            To connect this microscopic picture to transport, we consider how Bloch electrons move in weak external fields. On length and time scales large compared to the lattice spacing and inverse band gaps, the dynamics of an electron wave packet built from states in a single band $n$ can be described semiclassically by
            \begin{equation}
                \hbar \frac{\mathrm{d}\vec{k}}{\mathrm{d}t} = -e\left(
                    \vec{E} + \frac{\mathrm{d}\vec{r}}{\mathrm{d}t} \times \vec{B}
                \right)\,,
                \qquad
                \frac{\mathrm{d}\vec{r}}{\mathrm{d}t} = \frac{1}{\hbar}\,\nabla_{\vec{k}} E_n(\vec{k})\,,
                \label{eq:quasi:semiclassical}
            \end{equation}
            where $\mathrm{d}\vec{r}/\mathrm{d}t$ is the group velocity, given by the gradient of the band energy, and $\vec{E}$, $\vec{B}$ are external electric and magnetic fields. \cite{ashcroft_solid_2012}

            In a perfect, infinite crystal without interactions, Bloch electrons would evolve according to Eq.~\eqref{eq:quasi:semiclassical} and, for a uniform electric field, undergo Bloch oscillations rather than exhibit steady Ohmic transport. Real materials contain impurities, lattice defects, and phonons, and residual electron-electron interactions. These processes scatter carriers between Bloch states and thereby relax the current. \cite{bloch_uber_1929,zener_theory_1934}

            For the transport regimes relevant in this thesis it is sufficient to summarize all momentum-relaxing processes by a single energy-averaged momentum-relaxation time $t_\mathrm{qp}$. Between scattering events, electrons propagate as Bloch wave packets. Scattering events then randomize the momentum distribution and permit a stationary current. We now recast this picture in terms of quasiparticles and the effective parameters that connect band dynamics to measured conductance. \cite{ashcroft_solid_2012, gross_festkorperphysik_2014}

        %=========================================================
        \subsubsection*{Quasiparticles and Effective Single-Band Description}
        %=========================================================

            The independent-electron band picture treated interactions only at a mean-field level. In reality, the charged fermions in the Hamiltonian, which we will refer to as \emph{electrons}, interact with each other and with lattice vibrations. The low-energy excitations of a metal are therefore not bare electrons but \emph{quasiparticles}. In Landau's Fermi-liquid picture these quasiparticles are long-lived fermionic modes that carry charge $-e$ and spin $1/2$, but are dressed by their interaction with the surrounding Fermi sea. \cite{landau_theory_1956}

            Throughout this thesis we characterize normal-state quasiparticles in a narrow energy window around $E_\mathrm{F}$ by a small set of effective parameters: an effective mass $m^\ast$, a Fermi velocity $v_\mathrm{F}$, a constant normal-state density of states per unit volume $N_0 = N(E_\mathrm{F})$ including both spin directions, and a momentum-relaxation time $t_\mathrm{qp}$. The corresponding elastic mean free path is $\ell_\mathrm{qp}= v_\mathrm{F} t_\mathrm{qp}$. Here, $t_\mathrm{qp}$ is a momentum-relaxation time and need not coincide with the single-particle lifetime that governs spectral broadening.

            In an isotropic single-band approximation we approximate the dispersion by a parabolic form,
            \begin{equation}
                E(\vec{k}) \approx \frac{\hbar^2 \lvert \vec{k} \rvert^2}{2m^\ast}\,,
                \qquad
                E_\mathrm{F} = \frac{\hbar^2 k_\mathrm{F}^2}{2m^\ast}\,.
                \label{eq:quasi:parabolic}
            \end{equation}
            The effective mass generally differs from the bare electron mass $m$ and encodes both the band-structure curvature and renormalization due to interactions. The quasiparticle group velocity is then
            \begin{equation}
                \vec{v}(\vec{k})
                = \frac{1}{\hbar} \nabla_{\vec{k}} E(\vec{k})
                \approx \frac{\hbar\vec{k}}{m^\ast}\,,
                \label{eq:quasi:quasiparticle-velocity}
            \end{equation}
            and the Fermi velocity is $v_\mathrm{F} = \hbar k_\mathrm{F} / m^\ast$.

            The density of states per unit volume, including both spin directions, for a three-dimensional parabolic band reads
            \begin{equation}
                N(E) = \frac{1}{2\pi^2} \left( \frac{2m^\ast}{\hbar^2} \right)^{3/2} \sqrt{E}\,,
                \label{eq:quasi:n-dos}
            \end{equation}
            with $E$ measured from the band edge. Its value at the Fermi energy is $N_0 = N(E_\mathrm{F})$. Since $E_\mathrm{F}$ is of order eV whereas the relevant excitation energies are on the meV scale, $N(E)$ varies only weakly in the window of interest. It is therefore convenient to approximate it by a constant,
            \begin{equation}
                N(E) \approx N_0 = N(E_\mathrm{F})\,,
                \label{eq:quasi:n-dos-approx}
            \end{equation}
            This approximation is accurate in the energy range relevant for low-temperature transport. The parameters $m^\ast$, $v_\mathrm{F}$, $N_0$, and $t_\mathrm{qp}$ define the minimal quasiparticle description used throughout this thesis. They determine both spectral quantities and the normal-state transport scales, providing the input for the Drude model derived next. \cite{ashcroft_solid_2012, gross_festkorperphysik_2014}

        %=========================================================
        \subsubsection*{Drude Model from the Microscopic Picture}
        %=========================================================

            Within the semiclassical picture the simplest transport model is the Drude description of a homogeneous metal. We treat the quasiparticles introduced above as semiclassical carriers with charge $-e$, effective mass $m^\ast$, and density $n$. Random momentum-relaxing collisions are represented by the relaxation time $t_\mathrm{qp}$. In a weak, static electric field $\vec{E}$ the equation of motion for the average quasiparticle momentum reads
            \begin{equation}
                \frac{\mathrm{d}\langle \vec{p}\, \rangle}{\mathrm{d}t}
                = -e \vec{E} - \frac{\langle \vec{p}\, \rangle}{t_\mathrm{qp}}\,,
                \label{eq:quasi:drude-eom}
            \end{equation}
            where the second term is a simple relaxation-time approximation for the effect of scattering. In the steady state, $\mathrm{d}\langle \vec{p} \,\rangle / \mathrm{d}t = 0$ and $\langle \vec{p} \rangle = - e t_\mathrm{qp} \vec{E}$. The average drift velocity is therefore
            \begin{equation}
                \vec{v}_\mathrm{drift}
                = \frac{\langle \vec{p}\, \rangle}{m^\ast}
                = - \frac{e t_\mathrm{qp}}{m^\ast} \vec{E}\,.
                \label{eq:quasi:drift-velocity}
            \end{equation}
            Multiplying by the carrier density $n$ gives the current density
            \begin{equation}
                \vec{j} = - n e \vec{v}_\mathrm{drift}
                = \frac{n e^2 t_\mathrm{qp}}{m^\ast} \vec{E}\,,
                \label{eq:quasi:drude-current-density}
            \end{equation}
            which identifies the Drude conductivity
            \begin{equation}
                \sigma = \frac{n e^{2} t_\mathrm{qp}}{m^\ast}\,.
                \label{eq:quasi:drude-conductivity}
            \end{equation}
            The conductivity therefore translates the microscopic relaxation time into the slope of the normal-state \textit{I--V} characteristic. In this bulk limit the system is assumed to be spatially homogeneous, and spatial variations of the electric field and current density are slow on the scale of the microscopic scattering lengths, such that transport is fully captured by the local constitutive relation $\vec{j} = \sigma \vec{E}$.

            Integrating the local constitutive relation $\vec{j}=\sigma\vec{E}$ for a homogeneous wire of length $L$ and cross-sectional area $A$ yields the familiar macroscopic Ohm's law. In the conductance form that will be used throughout this thesis, one writes
            \begin{equation}
                I = G \ V\,,\qquad
                G = \sigma\,\frac{A}{L}\,.
                \label{eq:quasi:ohms-law}
            \end{equation}
            Equivalently, $V=IR$ with $R = 1/G$ and one may introduce the resistivity $\rho= 1/\sigma$ such that $R=\rho L/A$. Thus, a measured normal-state resistance combines the microscopic scattering rate with the sample geometry. This bulk relation provides the reference against which coherent mesoscopic transport is identified. \cite{ohm_galvanische_1827}

            The same relaxation time $t_\mathrm{qp}$ that enters $\sigma$ also sets the characteristic length and time scales of quasiparticle motion. Momentum randomization over this time defines an elastic mean free path
            \begin{equation}
                \ell_\mathrm{qp}= v_\mathrm{F} t_\mathrm{qp}\,,
                \label{eq:quasi:mean-free-path}
            \end{equation}
            where $v_\mathrm{F}$ is the Fermi velocity, and on longer scales the coarse-grained dynamics of the carrier density become diffusive with diffusion constant
            \begin{equation}
                D = \frac{v_\mathrm{F}^{2} t_\mathrm{qp}}{3} = \frac{v_\mathrm{F}\ell_\mathrm{qp}}{3}\,.
                \label{eq:quasi:diffusion-constant}
            \end{equation}
            The quantities $\sigma$, $\ell_\mathrm{qp}$, and $D$ summarize the influence of microscopic scattering in the bulk diffusive limit. The comparison of $\ell_\mathrm{qp}$ with the conductor size $L$ determines when this local description remains adequate. Once finite size and phase coherence matter, the transport problem must be organized by a hierarchy of length scales rather than by conductivity alone. \cite{drude_zur_1900}
                
    %=========================================================
    \newpage
    \subsection{Mesoscopic Quasiparticles}
    \label{subsec:quasi:meso}
    %=========================================================

        Mesoscopic transport occurs when the conductor size $L$ becomes comparable to microscopic transport length scales. Finite-size effects and quantum interference can then produce behavior that is not captured by the bulk Drude picture. The relevant hierarchy involves the elastic mean free path $\ell_\mathrm{qp}$, the phase-coherence length $\ell_\phi$, and the thermal length $\ell_T$. This hierarchy determines the applicable transport description.

        The phase-coherence length $\ell_\phi$ is set by inelastic and dephasing processes such as electron--electron and electron--phonon scattering. It quantifies how far a quasiparticle propagates while retaining a well-defined quantum phase. At finite temperature, thermal averaging further limits observable interference on the scale $\ell_T$.

        In macroscopic conductors, where $L \gg \ell_\mathrm{qp}$ and $L \gg \ell_\phi$, quasiparticles undergo many scattering events and lose phase coherence before traversing the sample. Transport is then well described by the bulk Drude picture. When $L \lesssim \ell_\phi$, phase coherence extends across the conductor and quantum interference between scattering paths becomes observable. The mean free path then distinguishes two coherent regimes. Transport is ballistic for $L \ll \ell_\mathrm{qp}$ and diffusive but phase coherent for $\ell_\mathrm{qp} \ll L \lesssim \ell_\phi$. Figure~\ref{fig:quasi:meso} summarizes this hierarchy. \cite{ashcroft_solid_2012, gross_festkorperphysik_2014}

        \begin{figure}
            \centering
            \subfigure[
                Classical Regime
                ]{\import{theory/schema/scattering}{classic.pdf_tex}}
            \hfill
            \subfigure[
                Diffusive Regime
                ]{\import{theory/schema/scattering}{diffusive.pdf_tex}}
            \hfill
            \subfigure[
                Ballistic Regime
                ]{\import{theory/schema/scattering}{ballistic.pdf_tex}}
            \caption{
            Transport regimes are set by conductor length $L$ relative to elastic mean free path $\ell_{\mathrm{qp}}$ and phase-coherence length $\ell_{\phi}$. (a)~Classical regime with $L \gg \ell_{\phi}, \ell_{\mathrm{qp}}$. (b)~Phase-coherent diffusive regime with $\ell_{\mathrm{qp}} \ll L \lesssim \ell_{\phi}$. (c)~Ballistic regime with $L \ll \ell_{\mathrm{qp}}$ and $L \lesssim \ell_{\phi}$.
            }
            % Legend: constriction \legend{seegrau120}, scattering centers \legend{seegrau100}, phase-coherent trajectory \legend{seeblau65}, phase-randomized trajectory \legend{seegrau65}, phase-coherence disk \legend{seeblau20}.
            \label{fig:quasi:meso}
        \end{figure}

        %=========================================================
        \subsubsection*{Diffusive Mesoscopic Transport}
        %=========================================================

            In the diffusive mesoscopic regime the elastic mean free path remains much shorter than the system size, $\ell_\mathrm{qp} \ll L$, so that quasiparticles undergo many momentum-randomizing scattering events while traversing the conductor. At the same time, phase coherence can persist over many such events, $L \lesssim \ell_\phi$, such that quantum interference between multiple scattering paths becomes observable.

            For diffusive motion the phase-coherence length is conveniently expressed in terms of the diffusion constant $D$ and a dephasing time $t_\phi$,
            \begin{equation}
                \ell_\phi = \sqrt{D t_\phi}\,.
                \label{eq:quasi:diffusive-coherence-length}
            \end{equation}
            At finite temperature, interference is additionally reduced by thermal averaging over an energy window $\sim k_\mathrm{B}T$. This introduces the thermal length
            \begin{equation}
                \ell_T = \sqrt{\hbar D / k_\mathrm{B}T}\,,
                \label{eq:quasi:thermal-length}
            \end{equation}
            such that phase-coherent effects are most pronounced when $L \lesssim \min(\ell_\phi,\,\ell_T)$.

            On average, transport still follows the classical drift-diffusion picture with conductance and diffusion constant as introduced in the Drude model above. Quantum mechanically, however, the coherent superposition of scattering amplitudes produces characteristic corrections to the Drude conductance, most prominently weak localization and universal conductance fluctuations. In low-dimensional or strongly disordered systems these interference effects can become strong and may ultimately lead to Anderson localization.
            \cite{bergmann_weak_1984, lee_universal_1985, lee_disordered_1985, anderson_absence_1958, schertel_magnetic-field_2019}

            In the present work the diffusive mesoscopic regime serves as a reference rather than the primary contact model. The atomic-scale constrictions studied here are shorter than the relevant scattering lengths, which leads to the ballistic channel description introduced next.

        %=========================================================
        \subsubsection*{Ballistic Transport and Scattering Formalism}
        %=========================================================
            
            \begin{wrapfigure}[14]{r}{0.4\textwidth}
                \captionsetup{format=plain}%
                \centering
                \vspace{-1em} % fine-tune vertical position
                \import{theory/schema/scattering}{scattering-matrix.pdf_tex}
                \vspace{-.25em} % fine-tune vertical position
                \caption{
                    Sketch of a two-terminal scatterer connecting ideal leads. Incoming and outgoing mode amplitudes $\vec{A}$, $\vec{B}$ are related by the scattering matrix $\mat{S}$.
                }
                \label{fig:quasi:scatter}
            \end{wrapfigure}
            In the ballistic regime the length of the constriction is shorter than the elastic mean free path, $L \ll \ell_\mathrm{qp}$, such that quasiparticles traverse the contact without momentum-randomizing scattering. If, in addition, the contact is shorter than the phase-coherence length, $L \lesssim \ell_\phi$, transport is phase coherent across the entire constriction. In the ballistic limit one may express the phase-coherence length in terms of the Fermi velocity and a dephasing time $t_\phi$,
            \begin{equation}
                \ell_\phi = v_\mathrm{F} \, t_\phi\,.
                \label{eq:quasi:ballistic-coherence-length}
            \end{equation}

            Ballistic transport is most naturally formulated as an elastic scattering problem. A finite region connects ideal leads in which the electronic modes propagate without dissipation. We denote the complex incoming mode amplitudes in the left and right leads by $A_1$ and $A_2$, and the corresponding outgoing amplitudes by $B_1$ and $B_2$, as shown in Fig.~\ref{fig:quasi:scatter}.
            
            Linear scattering then implies a matrix relation which we write explicitly as
            \begin{equation}
                \begin{pmatrix} B_1 \\ B_2 \end{pmatrix} = \mat{S}\,\begin{pmatrix} A_1 \\ A_2 \end{pmatrix}\,.
                \label{eq:quasi:scatter-1}
            \end{equation}

            In the general multi-mode case, $A_{1,2}$ and $B_{1,2}$ become vectors collecting all propagating modes in the corresponding lead. The scattering matrix decomposes into reflection blocks $\mat{r}$ and $\mat{r}'$ and transmission blocks $\mat{t}$ and $\mat{t}'$,
            \begin{equation}
                \mat{S}
                = \begin{pmatrix} \mat{r} & \mat{t}' \\ \mat{t} & \mat{r}' \end{pmatrix}\,.
                \label{eq:quasi:scatter-2}
            \end{equation}
            For purely elastic normal-state scattering, $\mat{S}$ is unitary because probability current is conserved.
            
            The central quantities governing transport are then the transmission eigenvalues
            \begin{equation}
                \{\tau_i\} = \operatorname{eig}\!\left(\mat{t}^\dagger\,\mat{t}\right)\,,\quad 
                \tau_i\in[0,1]\,.
                \label{eq:quasi:tau_i}
            \end{equation}
            Choosing a basis of transmission eigenchannels diagonalizes $\mat{t}^\dagger\mat{t}$. This basis is unique up to unitary rotations within degenerate subspaces.

            In linear response, the normal-state conductance is then given by the Landauer formula,
            \begin{equation}
                G_\mathrm{N} = G_0 \sum_{i=1}^N \tau_i\,,
                \quad
                G_0 = \frac{2e^2}{h}\,,
                \label{eq:quasi:landauer}
            \end{equation}
            where $\tau_i$ are the transmission eigenvalues and $G_0$ is the conductance quantum\footnote{
                Throughout this thesis the conductance quantum is defined as $G_0 \equiv 2e^2/h = 77.48\,\mu\mathrm{S}$. This value includes spin degeneracy and corresponds to the conductance of a fully transmitting normal-state channel. The corresponding resistance quantum is $R_0 = h/(2e^2) = 12.9\,\mathrm{k\Omega}$.
            }.

            The normal-state conductance measures only the sum of the transmission eigenvalues. It does not determine their individual values. The full channel set $\{\tau_i\}$, often referred to as the mesoscopic pin code of the contact, becomes observable through the channel-dependent superconducting \textit{I--V} and \textit{dI--dV} characteristics discussed later. This distinction connects the Landauer formula to the experimental channel extraction. \cite{landauer_electrical_1970,buttiker_four-terminal_1986, beenakker_random-matrix_1997}

        %=========================================================
        \subsubsection*{Atomic-Scale Contacts}
        %=========================================================

            As a didactic reference for the channel picture, it is useful to recall the textbook case of a split-gate quantum point contact (QPC) in a high-mobility two-dimensional electron gas. There, the width of a short ballistic constriction is tuned electrostatically, and the conductance increases in approximately quantized steps as additional transverse modes become available. In the Landauer picture each newly populated mode contributes roughly one conductance quantum $G_0$ when its transmission approaches unity. The resulting staircase $G(V_\mathrm{G})$ is shown in Fig.~\ref{fig:quasi:atomic-contact}(a). \cite{van_wees_quantized_1988, buttiker_four-terminal_1986, beenakker_random-matrix_1997}

            While semiconductor QPCs provide a particularly clean and tunable example of conductance quantization, the focus of this work lies on metallic atomic contacts. In these contacts the number, symmetry, and transparency of the transmission channels are dictated by atomic-scale structure and chemistry rather than by lithographic geometry.

            \begin{figure}
                \centering
                \subfigure[
                    QPC conductance $G$ vs. gate voltage $V_\mathrm{G}$, showing quantized plateaus.
                    \cite{van_wees_quantized_1988}
                    ]{\import{theory/basics}{GV.pgf}}
                \hfill
                \subfigure[
                    Aluminum atomic contact opening trace, conductance $G$ vs. displacement $\Delta x$.
                    \cite{scheer_conduction_1997}
                    ]{\import{theory/basics}{Gx.pgf}}
                \hfill
                \subfigure[
                    Conductance histogram $P(G)$ for aluminum atomic contacts.
                    \cite{yanson_histograms_1997}
                    ]{\import{theory/basics}{PG.pgf}}

                \caption{
                    Representative conductance quantization in mesoscopic constrictions and atomic contacts.
                (a) Split-gate quantum point contact: $G(V_\mathrm{G})$.
                (b) Aluminum atomic contact: opening trace $G(\Delta x)$.
                (c) Conductance histogram $P(G)$ accumulated over many breaking cycles.
                }
                \label{fig:quasi:atomic-contact}
            \end{figure}

            Metallic atomic contacts provide a paradigmatic realization of ballistic, phase-coherent transport. The constriction consists of only one or a few atoms, with a characteristic length of order an interatomic spacing. This scale is far shorter than the elastic mean free path and, at low temperature, also much shorter than the phase-coherence length, such that typically $L \ll \ell_\mathrm{qp}$ and $L \ll \ell_\phi$. Experimentally, the formation and rupture of such contacts are tracked by recording the conductance while mechanically separating the electrodes, producing an opening trace as in Fig.~\ref{fig:quasi:atomic-contact}(b). \cite{scheer_conduction_1997, agrait_quantum_2003}

            The transmission eigenchannels of an atomic contact are set by the valence orbitals of the atoms forming the narrowest part of the constriction and by their hybridization with the electrodes. For aluminum, whose conduction states have predominantly $sp$ character, this typically results in three dominant transport channels: one mainly $sp_z$-like channel and two transverse $p$-like channels. Depending on atomic geometry and disorder, additional weaker channels may appear, but the transport is usually governed by a small number of highly transparent modes. A complementary statistical view is provided by conductance histograms, Fig.~\ref{fig:quasi:atomic-contact}(c), where peaks reflect preferred atomic configurations and their associated typical channel sets. \cite{yanson_histograms_1997,scheer_signature_1998, agrait_quantum_2003}

            Experimentally, the set of transmission probabilities $\{\tau_i\}$ can be extracted from superconducting transport measurements using multiple Andreev reflection and related techniques. Throughout this thesis we therefore characterize each atomic contact by its mesoscopic pin code. These transmission eigenvalues serve as input for the microscopic description of superconducting transport in the following sections. \cite{scheer_conduction_1997,scheer_signature_1998}

            Together, opening traces and conductance histograms show how atomic configurations select reproducible conductance ranges. They motivate the channel description used throughout this thesis, while superconducting subgap spectroscopy supplies the additional information required to resolve the individual $\tau_i$. The remaining step is to specify the material scales that place aluminum contacts in this ballistic superconducting regime.

    In summary, the comparison of $L$ with $\ell_\mathrm{qp}$, $\ell_\phi$, and $\ell_T$ determines whether transport is described by a local Drude conductivity or by a phase-coherent scattering problem. In the coherent regime, the transmission eigenvalues $\{\tau_i\}$ connect the contact structure to normal-state conductance and superconducting subgap observables. Section~\ref{sec:platform} now specializes this general framework to aluminum and introduces the voltage and phase conventions used for driven junctions.


